\documentclass[11pt]{article}
\usepackage[utf8]{inputenc}
\usepackage[T1]{fontenc}
\usepackage{geometry}
\geometry{margin=1in}
\usepackage{microtype}
\usepackage{amsmath, amssymb, amsthm, mathtools}
\usepackage{array}
\usepackage{booktabs}
\usepackage{enumitem}
\usepackage{xurl}
\usepackage[hidelinks]{hyperref}

\setlist{itemsep=0.25em, topsep=0.4em}
\emergencystretch=2em

\title{The General Theory of Finite Descent Authority}
\author{Jeremy H. Carroll}
\date{May 2026}

\newtheorem{theorem}{Theorem}[section]
\newtheorem{lemma}[theorem]{Lemma}
\newtheorem{definition}[theorem]{Definition}
\newtheorem{axiom}[theorem]{Axiom}
\newtheorem{claim}[theorem]{Claim}
\newtheorem{proposition}[theorem]{Proposition}

\newcommand{\boxedtext}[1]{%
\begin{center}
\fbox{\begin{minipage}{0.92\linewidth}\centering #1\end{minipage}}
\end{center}}

\begin{document}

\maketitle

\section*{Publication Main}

\begin{quote}
All files are available open-source on \url{https://github.com/dbccpage/general_theory_of_finite_obstruction_calculus}
\end{quote}

\begin{quote}
This file is the publication-facing spine extracted from the working tome. Full proofs are in \path{2_GTFDA_Proof_Archive.md}.

LLM, Kaggle, domain, quantum, and observer-selection material is in \path{4_GTFDA_Tertiary_Workbench.md}.
\end{quote}

\section*{Relevant High-Level Papers}

\begin{enumerate}[label={[\arabic*]}]
    \item J. H. Carroll. \textit{General Theory of Finite Obstruction Calculus} 2026. Zenodo. \url{https://zenodo.org/uploads/20193705}
    \item J. H. Carroll. \textit{General Theory of Layered Authority} 2026. Zenodo.
\end{enumerate}

\section*{Reader's Guide}

\textbf{Finite Descent Authority (FDA)} is the doctrine that a proposal earns operational authority only by entering a finite certified regime, being verified inside that regime, and terminalizing into a declared authority-bearing status. The main paper keeps only the core definitions and theorem statements needed for the doctrine. The proof archive preserves the full proof text. The tertiary workbench keeps applications, implementation-specific ledgers, and speculative extensions out of the publication spine.

\section*{The Core Normal Form}

The stable authority identity is:

\[
\boxed{
\mathsf{Auth}(p)=
\tau_s\bigl(V_s(P_s,\eta_s(p))\bigr),
\qquad
s=\rho(p),
\quad
P_s=D(s).
}
\]

Equivalently, if \(V_s(P_s,\eta_s(p))=(R_p,O_p,\varepsilon_s)\), then:

\[
\boxed{
\mathsf{Auth}(p)=\tau_s(R_p,O_p,\varepsilon_s).
}
\]

The compressed notation \(\mathsf{Auth}=\tau\circ V\circ D\) suppresses routing and sector-local realization. The theorem-grade reading always includes \(p\mapsto s=\rho(p)\), \(\eta_s(p)\), \(P_s=D(s)\), verification, and terminalization.

\begin{abstract}
A system does not acquire authority by producing a representation, a prediction, a fluent explanation, or a global geometry. It acquires authority only when the claim descends into a finite certified regime with declared carriers, operators, verifiers, terminal statuses, and refusal conditions.

The governing identity is:

\[
\boxed{
\mathsf{Auth}
=
\tau\circ V\circ D
}
\]

The refusal law is:

\[
\boxed{
\neg D\Rightarrow\neg\mathsf{Auth}.
}
\]

Finite obstruction calculus is the local engine of this theory. It analyzes defects, repairs, quotient obstructions, witnesses, and licensed lifts inside a descended finite packet. Finite descent authority is the broader doctrine determining when any high-level proposal earns operational authority.
\end{abstract}

\section*{Non-Claims}

This theory does not claim that every truth is finite. It does not claim that every domain already has a known descent map. It does not claim that global theories, geometric structures, or LLMs are useless. It claims that operational authority, as defined here, requires explicit finite descent, verification, terminal semantics, and refusal conditions.

\section{Notation and Packet Levels}

This paper uses two distinct maps. The router is:
\[
\rho:\mathcal P\to\mathcal G.
\]
It sends a proposal to a sector, family, shape, regime, or domain-local class. It has no authority.

Sector-local realization is written:
\[
\eta_s:\mathcal P_s\to\mathcal R_s.
\]
It encodes a proposal into the finite carrier, observer, operator, diagnostic, and replay data of sector \(s\). The theorem-grade reading is always:
\[
s=\rho(p),
\qquad
\eta_s(p)\downarrow.
\]

The finite descent map is:
\[
D:\mathcal G\to\mathcal C_{\mathrm{fin}}.
\]
For a sector \(s\), it returns the abstract authority packet:
\[
D(s)=P_s=(K_s,E_s,V_s,\tau_s,\Lambda_s).
\]

The expanded obstruction-engine tuple:
\[
(P,E,\mathcal A,\Phi,\Gamma,\Xi,\mathcal{NA},\mathcal{RA},R_{cl},C_{\mathrm{rollback}},\tau)
\]
is not a replacement for \(P_s=(K_s,E_s,V_s,\tau_s,\Lambda_s)\). It is an implementation refinement:
\[
(K_s,E_s,V_s,\tau_s,\Lambda_s) \leadsto \operatorname{Impl}(P_s),
\]
where:
\begin{itemize}
    \item \(K_s\) contains the finite carrier, observer/profile data, and representation state;
    \item \(E_s\) contains admissible operators and evaluator diagnostics such as \(\Phi,\Gamma, \Xi, \mathcal{NA}, \mathcal{RA}, R_{cl}\);
    \item \(V_s\) contains replay checks, certificates, rollback data, and audit evidence;
    \item \(\tau_s\) contains terminal semantics;
    \item \(\Lambda_s\) contains lift, refusal, backtrack, budget, and no-silent-mutation discipline.
\end{itemize}

Thus the stable authority identity is:
\[
\boxed{
\mathsf{Auth}(p)=
\tau_s\bigl(V_s(P_s,\eta_s(p))\bigr),
\qquad
s=\rho(p),
\quad
P_s=D(s).
}
\]
When \(V_s(P_s,\eta_s(p))=(R_p,O_p,\varepsilon_s)\), this becomes:
\[
\boxed{
\mathsf{Auth}(p)=\tau_s(R_p,O_p,\varepsilon_s).
}
\]

When the paper writes the compressed composition \(\tau\circ V\circ D\circ\rho\), the sector-local realization \(\eta_s\) is suppressed. The expanded reading is always:
\[
p\mapsto s=\rho(p) \mapsto P_s=D(s) \mapsto V_s(P_s,\eta_s(p)) \mapsto \tau_s.
\]

\part{Core Theorems}

\section{FDA-1 — Proposal Non-Authority Theorem}

\begin{claim}
\[
\boxed{
p\in\mathcal P \not\Rightarrow \mathsf{Auth}(p)
}
\]
A proposal object has no operational authority merely by existing, regardless of whether it came from an LLM, a theory, an embedding, a proof sketch, a policy model, a database query, a program, or a global geometric construction.
\end{claim}

\paragraph{Setup}
Let \(\mathcal P\) be the class of proposal objects. A proposal \(p\in\mathcal P\) is only an offered candidate. It may be a transition, proof step, repair, lift, rewrite, policy, model update, program transformation, or reasoning trace.

Authority is not defined directly on \(\mathcal P\). It is produced only through the finite authority pipeline:
\[
\mathcal P \xrightarrow{\rho} \mathcal G \xrightarrow{D} \mathcal C_{\mathrm{fin}} \xrightarrow{V} \mathsf{Fib} \xrightarrow{\tau} \mathsf{Term}.
\]
Where:
\begin{itemize}
    \item \(s=\rho(p)\in\mathcal G\) is the routed sector for \(p\);
    \item \(D(s)=P_s\in\mathcal C_{\mathrm{fin}}\) is the finite descent packet for that sector, if available;
    \item \(\eta_s(p)\) is the sector-local finite realization of \(p\);
    \item \(V_s(D(s),\eta_s(p))\in\mathsf{Fib}\) is the verified fiber;
    \item \(\tau\) terminalizes the verified fiber;
    \item \(\mathsf{Auth}(p)\) is an authority-bearing terminal output.
\end{itemize}

Thus:
\[
\mathsf{Auth}(p) \iff \tau_s(V_s(D(s),\eta_s(p)))\in\mathsf{Term}_{\mathsf{auth}}, \qquad s=\rho(p).
\]
Authority is a terminal status, not a property of proposal existence.

\begin{theorem}
\[
\boxed{
p\in\mathcal P \not\Rightarrow \mathsf{Auth}(p)
}
\]
\end{theorem}

\section{FDA-2 — Finite Descent Necessity Theorem}

\begin{claim}
\[
\boxed{\neg D \Rightarrow \neg \mathsf{Auth}}
\]
No finite descent packet, no finite operational authority.
\end{claim}

\paragraph{Definitions}
Let \(p\) be a proposed transition, repair, lift, rewrite, policy, proof step, or operator promotion. Let \(s=\rho(p)\) denote its routed sector. Let \(\eta_s(p)\) denote the finite realization of \(p\) into the sector-local obstruction engine: carrier, observer, operators, diagnostics, rollback data, and evaluator state. Let \(D(s)\downarrow\) mean that the routed sector has a finite descent packet.

A finite descent packet contains enough replayable data to decide one of the authorized operational routes: \(\mathsf{accept}, \mathsf{reject}, \mathsf{lift}, \mathsf{backtrack}, \mathsf{hold}, \mathsf{unknown}\). At the abstract level:
\[
D(s)=P_s = (K_s,E_s,V_s,\tau_s,\Lambda_s).
\]
At the obstruction-engine implementation level this packet may refine to:
\[
\operatorname{Impl}(P_s) = (P,E,\mathcal A,\Phi,\Gamma,\Xi,\mathcal{NA},\mathcal{RA},R_{cl},C_{\mathrm{rollback}},\tau_s).
\]

Define operational authority as a partial predicate \(\mathsf{Auth}(p)\), defined only when \(D(s)\downarrow\), where \(s=\rho(p)\). Equivalently, \(\operatorname{dom}(\mathsf{Auth}) = \{p : s=\rho(p),\ D(s)\downarrow\}\). Thus, \(\mathsf{Auth}(p) \Rightarrow D(s)\downarrow\).

\begin{theorem}
\[
\boxed{
\neg D(s)\downarrow \Rightarrow \neg \mathsf{Auth}(p)
}
\qquad s=\rho(p).
\]
\end{theorem}

\section{FDA-3 — Authority Factorization Theorem}

\begin{claim}
\[
\boxed{\mathsf{Auth}=\tau\circ V\circ D}
\]
More precisely: \(\mathsf{Auth}(p)=\tau_s(R_p,O_p,\varepsilon_s)\), where \(V_s(D(s),\eta_s(p))=(R_p,O_p,\varepsilon_s)\).
\end{claim}

\paragraph{Objects}
Let \(p\) be a proposed transition and \(s=\rho(p)\) be its routed sector. Let \(\eta_s(p)\) be the finite realization of \(p\). Let \(D(s) = P_s = (K_s,E_s,V_s,\tau_s,\Lambda_s)\) be the finite descent packet.

\paragraph{Verification Map}
Define the verification map \(V_s(D(s),\eta_s(p)) = (R_p,O_p,\varepsilon_s)\), where \(R_p\) is route data, \(O_p\) is obstruction data, and \(\varepsilon_s\) is the certified error/exception boundary. Possible route states include: \(\mathsf{accept}, \mathsf{reject}, \mathsf{lift}, \mathsf{backtrack}, \mathsf{hold}, \mathsf{unknown}\).

\paragraph{Terminalization Map}
Define \(\tau_s\) as the terminalization map for sector \(s\). It consumes the verified fiber and returns an authority-bearing terminal state: \(\tau_s(R_p,O_p,\varepsilon_s) \in \mathsf{Terminal}_s\). Examples: \(\mathsf{AcceptCertified}, \mathsf{RejectCertified}, \mathsf{LicensedLift}, \mathsf{CertifiedBacktrack}, \mathsf{Hold}, \mathsf{BudgetUnknown}, \mathsf{TypedRefusal}\).

\begin{theorem}
\[
\boxed{
\mathsf{Auth}(p) = \tau_s(V_s(D(s),\eta_s(p)))
}
\]
Equivalently, if \(V_s(D(s),\eta_s(p)) = (R_p,O_p,\varepsilon_s)\), then:
\[
\boxed{
\mathsf{Auth}(p) = \tau_s(R_p,O_p,\varepsilon_s).
}
\]
\end{theorem}

\section{FDA-4 — Descent Packet Soundness}

\begin{claim}
A descent packet \(P_s=(K_s,E_s,V_s,\tau_s,\Lambda_s)\) is authority-bearing only if each component is declared and replayable:
\[
\boxed{
\mathsf{AuthPacket}(P_s) \Rightarrow \mathsf{DeclaredReplayable}(K_s,E_s,V_s,\tau_s,\Lambda_s).
}
\]
\end{claim}

\begin{theorem}
\[
\boxed{
\mathsf{AuthPacket}(P_s) \Rightarrow K_s\downarrow \wedge E_s\downarrow \wedge V_s\downarrow \wedge \tau_s\downarrow \wedge \Lambda_s\downarrow
}
\]
where \(X\downarrow\) means the component is declared, finite or finitely auditable, and replayable in the packet.
\end{theorem}

\section{FDA-5 — Terminal Authority Theorem}

\begin{claim}
The authority-bearing object is not the answer alone, but the terminal pair:
\[
\boxed{
(\mathsf{Term},\ \text{derived output if any})
}
\]
An output without its terminal status has no reasoning authority.
\end{claim}

\begin{theorem}
\[
\boxed{
\mathsf{Auth}(p) \Rightarrow \exists \mathsf{Term}_p,\ y_p^{?} \text{ such that } \mathsf{AuthObj}(p)=(\mathsf{Term}_p,y_p^{?})
}
\]
and:
\[
\boxed{
y_p\not\Rightarrow \mathsf{Auth}(p)
}
\]
An answer alone is not authority.
\end{theorem}

\part{Relation to Finite Obstruction Calculus}

\section{FDA/GT-1 — Obstruction Engine Inclusion Theorem}

\begin{claim}
\[
\boxed{ \mathsf{FiniteObstructionCalculus} \subset \mathsf{FiniteDescentAuthority} }
\]
More precisely: every finite obstruction-calculus engine instance embeds into a finite descent authority packet \(P=(K,E,V,\tau,\Lambda)\) where obstruction calculus supplies the internal defect, repair, quotient, witness, and lift machinery, but authority is granted only after terminalization.
\end{claim}

\paragraph{Embedding Into FDA Packet}
Define a map
\[
\iota:\mathsf{FiniteObstructionCalculus}\to\mathsf{FiniteDescentAuthority}
\]
by sending an obstruction-calculus engine \(\mathcal M\) to the descent packet
\[
\iota(\mathcal M)=P_{\mathcal M}=(K,E,V,\tau,\Lambda).
\]
\begin{itemize}
    \item \(K_{\mathcal M} = (\mathcal T,\mathcal C,\rho,C^\bullet,\omega,p)\) supplies the finite obstruction carrier.
    \item \(E_{\mathcal M} = (\Phi,\Gamma,Q^1,H^1,\mathsf{UniqueRule},\mathsf{Survivor},\mathsf{Open},\mathcal A)\) supplies the obstruction evaluator.
    \item \(V_{\mathcal M}\) supplies certificate replay for \(d_1d_0=0\), \(\Phi\), \(\Gamma\), and dual witness validity.
    \item \(\tau_{\mathcal M}\) turns verified obstruction data into authority-bearing terminal status.
    \item \(\Lambda_{\mathcal M}\) supplies lift/refusal discipline based on context-visible non-congruence.
\end{itemize}

\begin{theorem}
\[
\boxed{ \mathsf{FiniteObstructionCalculus} \subset \mathsf{FiniteDescentAuthority} }
\]
\end{theorem}

\section{FDA/GT-2 — Verification as Exactness Theorem}

\begin{claim}
\[
\boxed{ V_s(r,x)=1 \Rightarrow [\omega_{r,x}]=0\in Q^1_s }
\]
under the packet's declared obstruction encoding.
\end{claim}

\begin{theorem}
This connects replay verification to quotient exactness. If \(V_s(r,x)=1\), then replay produces an exact repair certificate: \(\exists \alpha\in C^0_s\) s.t. \(\omega_{r,x}=d_{0,s}\alpha\).
\end{theorem}

\section{FDA/GT-3 — Failure Trichotomy Theorem}

\begin{claim}
\[
\boxed{ V_s(r,x)=0 \Rightarrow [\omega_{r,x}]\ne0 \vee d_{1,s}\omega_{r,x}\ne0 \vee \neg\operatorname{adm}_{\Lambda_s}(r,x) }
\]
Failed verification must be one of: quotient obstruction, closure failure, or admissibility failure.
\end{claim}

\begin{theorem}
Define \(V_s(r,x)=1\) iff \([\omega_{r,x}]=0 \wedge d_{1,s}\omega_{r,x}=0 \wedge \operatorname{adm}_{\Lambda_s}(r,x)\). Failed verification implies the negation of at least one check.
\end{theorem}

\section{FDA/GT-4 — Lift Authority Theorem}

\begin{claim}
A lift \(P_s\to P_{s^+}\) is authority-preserving only if it satisfies: \(\iota_s \text{ mono}\), \(\pi_s\circ\iota_s=\operatorname{id}\), \(\mu_{s^+}<\mu_s\), plus no-silent-mutation.
\end{claim}

\begin{theorem}
\[
\boxed{
\mathsf{AuthorityPreservingLift}(P_s\to P_{s^+}) \Rightarrow \iota_s\ \mathrm{mono} \wedge \pi_s\circ\iota_s=\operatorname{id} \wedge \mu_{s^+}<\mu_s \wedge \mathsf{NoSilentMutation}
}
\]
\end{theorem}

\part{Publication Boundary}

The following material has been intentionally excluded from the publication main and moved to the tertiary workbench: LLM-specific corollaries, Kaggle trace geometry, domain-specific transfer, quantum packet examples, and implementation scaffolding.

\section*{Final Thesis}

\boxedtext{A proposal is input. A descent packet is machinery. A verifier is evidence. Terminalization is authority.}

\boxedtext{No finite descent packet, no operational authority.}

\end{document}
